\documentclass[11pt]{article}
\usepackage[T1]{fontenc}
\usepackage{lmodern}
\usepackage[margin=1in]{geometry}
\usepackage{amsmath,amssymb,amsthm,mathtools,booktabs,array,longtable}
\usepackage[colorlinks=true,linkcolor=blue,citecolor=blue,urlcolor=blue,breaklinks=true]{hyperref}
\usepackage{xurl}
\newtheorem{theorem}{Theorem}
\newtheorem{corollary}{Corollary}
\newcommand{\C}{\mathbb C}
\newcommand{\OB}{\operatorname{OB}}
\newcommand{\adj}{\operatorname{adj}}
\newcommand{\GL}{\operatorname{GL}}
\title{Proof of Critch--Morton Conjecture 8\protect\\
\large (SIGMA 10 (2014), 095, Section 2, p.\ 6):\protect\\
Reconstruction from any consecutive three-site improper marginal}
\author{Haobo Ma\\\small ChronoAI Pte Ltd, Singapore
\and Wenlin Zhang\\\small National University of Singapore\\\small ChronoAI Pte Ltd, Singapore}
\date{November 3, 2026}
\begin{document}
\maketitle
\begin{abstract}
We prove Critch--Morton Conjecture 8 (SIGMA 10 (2014), 095, Section 2, p.\ 6):
every consecutive three-site improper marginalization is birational
for open-boundary, translation-invariant, binary bond-two matrix product states
of length $N\geq4$.
We give an explicit rational inverse over $\C$, including a nonempty principal
open domain and both inverse identities on the affine Zariski closures.
A four-by-four linear system reconstructs a three-site representation; inverse
powers of the sum of its two matrices then restore the omitted sites.
Consequences include seven-dimensional rational models, reconstruction from
any longer consecutive block with arbitrary fixed nonzero site-dependent
coherent contractions, and explicit equations characterizing the reconstruction
chart. Three retained amplitude indices are necessary for generic recovery.
\end{abstract}

\section{The affine statement}
For $A_0,A_1\in\C^{2\times2}$, a row $\alpha\in\C^{1\times2}$,
and a column $\beta\in\C^{2\times1}$, write
\begin{equation}\label{eq:state}
 \psi_w=\alpha A_{w_1}\cdots A_{w_N}\beta,
 \qquad w\in\{0,1\}^N.
\end{equation}
Let $X_N=\overline{\OB(2,2,N)}$ be the affine Zariski closure of these
amplitude tensors. This is the model of Critch and Morton~\cite[Section 2,
equation (2)]{CM}. For $0\leq a\leq N-3$, put $b=N-3-a$ and define
\begin{equation}\label{eq:marginal}
 (\mu_a\psi)_{ijk}=\sum_{u\in\{0,1\}^a}\sum_{v\in\{0,1\}^b}\psi_{uijkv}.
\end{equation}
These are sums of complex amplitudes, not partial traces of density matrices.
With $S=A_0+A_1$, distributivity gives the useful identity
\begin{equation}\label{eq:outside}
 (\mu_a\psi)_{ijk}=\alpha S^a A_iA_jA_k S^b\beta.
\end{equation}
In particular $\mu_a:X_N\to X_3$ is a polynomial map: it maps the parametrized
model to the three-site model, and polynomial equations extend to its closure.

\begin{theorem}\label{thm:main}
For every $N\geq4$ and $0\leq a\leq N-3$, the map $\mu_a:X_N\to X_3$
is birational. More precisely, there is a polynomial $g$ in the eight
three-site coordinates such that $U=X_3\cap\{g\ne0\}$ is nonempty, and a
rational map $F_{N,a}:U\to X_N$ is inverse to $\mu_a$ on
$X_N\cap\{g\circ\mu_a\ne0\}$.
The vanishing ideals of $X_N$ and $X_3$ are prime.
\end{theorem}

The analogue for binary hidden Markov models appears in
Critch~\cite[Theorems 4.1 and 5.11]{HMM}. The free boundary vectors in
\eqref{eq:state} require the two-sided correction in \eqref{eq:outside}; a
stochastic normalization is not assumed here. Czapli\'nski, Micha\l{}ek and
Seynnaeve~\cite[Section 3.4, Theorem 3.30 of the arXiv version]{CMS}
settle the separate periodic identifiability Conjecture 12 of~\cite{CM}.
More general boundary-matrix canonical forms are developed
in~\cite[Sections 3--5]{FLAC}; our argument concerns the particular
three-site map in \eqref{eq:marginal}. For periodic linear equations and
their invariant-theoretic structure, see~\cite{LS}.

\section{An explicit inverse and its proof}
For $t=(t_{ijk})\in\C^8$, define the slices $T_j=(t_{ijk})_{i,k=0,1}$ and
\begin{equation}\label{eq:kernel}
 K(t)=\begin{pmatrix}
 t_{000}&t_{001}&-t_{000}&-t_{100}\\
 t_{010}&t_{011}&-t_{001}&-t_{101}\\
 t_{100}&t_{101}&-t_{010}&-t_{110}\\
 t_{110}&t_{111}&-t_{011}&-t_{111}
 \end{pmatrix}.
\end{equation}
Let $H$ be its upper-left three-by-three submatrix, $d=\det H$, and
let $c_j$ be the determinant obtained by replacing column $j$ of $H$
by $-K_{\{0,1,2\},3}$, for $j=0,1,2$. Set
\begin{align}
 \widetilde\ell&=(c_2,d),& \widetilde z&=(c_0,c_1)^T,&
 W_{ik}&=\sum_{j=0}^1\widetilde\ell_jt_{jik},\label{eq:raw}\\
 N_j&=T_j\adj W,&J&=N_0+N_1,&\widetilde r&=W\widetilde z.\nonumber
\end{align}
All these quantities are polynomial in $t$. Our domain is specified by
\begin{equation}\label{eq:g}
 g=d\det W\det J\det T_0\,
 (t_{000}t_{011}-t_{001}t_{010})\ne0.
\end{equation}
On this domain put
\begin{equation}\label{eq:reconstruction}
 \ell=\widetilde\ell/d,\quad z=\widetilde z/d,\quad B=W/d,
 \quad M_j=T_jB^{-1},\quad r=Bz,\quad C=M_0+M_1.
\end{equation}
Then $B$ and $C$ are invertible, since
$M_j=dN_j/\det W$ and $C=dJ/\det W$. The inverse in Theorem~\ref{thm:main} is
\begin{equation}\label{eq:inverse}
 F_{N,a}(t)_w=\ell C^{-a}M_{w_1}\cdots M_{w_N}C^{-b}r,
 \qquad b=N-3-a.
\end{equation}

\begin{proof}[Proof of Theorem~\ref{thm:main}]
\emph{The quartic relation.}
First take $t_{ijk}=lA_iA_jA_kr$. Define
$P_{i,*}=lA_i$ and $Q_{*,k}=A_kr$, so $T_j=PA_jQ$.
On the open set where $P,Q$ are invertible, put
$e=lP^{-1}$ and $v=Q^{-1}r$. For each $i,j$ one has
\begin{equation}\label{eq:relation}
 \sum_k t_{ijk}v_k=lA_iA_jr=\sum_k e_kt_{kij}.
\end{equation}
Consequently $K(t)(v_0,v_1,e_0,e_1)^T=0$. This vector is nonzero since
$Q$ invertible implies $r\ne0$. Thus $\det K(t)=0$ on that parameter open
set. It is nonempty: $A_0=I$, $A_1=\left(\begin{smallmatrix}0&1\\1&0\end{smallmatrix}\right)$,
$l=(1,0)$, $r=(1,0)^T$ give $P=Q=I$. The parameter space is affine
and irreducible, so the polynomial $\det K(t)$ vanishes for all parameters,
and therefore on $X_3$.

\emph{Three-site reconstruction.}
Cramer's rule applied to the first three rows of $K$ gives
\begin{equation}\label{eq:cramer}
 K(t)(c_0,c_1,c_2,d)^T=(0,0,0,\det K(t))^T.
\end{equation}
The last component follows by expanding the determinant in its last row.
For $t\in U$, division by $d$ gives
\begin{equation}\label{eq:compatibility}
 \sum_k t_{ijk}z_k=\sum_k\ell_kt_{kij}=B_{ij}.
\end{equation}
The definition of $B$ also gives $\ell T_i=B_{i,*}$. Hence
\begin{equation}\label{eq:normalized}
 \ell M_i=e_i^T,\qquad M_kr=T_kz=B_{*,k},\qquad
 \ell M_iM_jM_kr=(M_jB)_{ik}=t_{ijk}.
\end{equation}
Here $e_i$ is the $i$th coordinate column. Thus the reconstructed matrices
and boundaries give exactly $t$ as a three-site tensor.

\emph{A right inverse for every length and position.}
Equation~\eqref{eq:inverse} is itself of the form~\eqref{eq:state}; it
therefore belongs to the parametrized $N$-site model. Summing over the
omitted sites multiplies its boundaries by $C^a$ and $C^b$.
These powers cancel their inverses, and \eqref{eq:normalized} yields
$\mu_aF_{N,a}(t)=t$ for every $t\in U$.

\emph{Uniqueness on parametrized tensors.}
Let $t=\mu_a\psi\in U$ with $\psi$ as in \eqref{eq:state}.
Use $l=\alpha S^a$, $r_0=S^b\beta$, and form $P,Q$ as above with
$r_0$ in place of $r$. Since $\det T_0\ne0$ and $T_0=PA_0Q$,
both $P,Q$ are invertible. Define $e=lP^{-1}$ and $v=Q^{-1}r_0$.
The first three rows of $K$ are independent, because $d\ne0$.
Their kernel, and hence the kernel of $K$, is the line spanned by
$(c_0,c_1,c_2,d)^T$. The nonzero vector $(v,e)$ lies on this line.
Thus $s=e_1\ne0$, $\ell=s^{-1}e$, and $z=s^{-1}v$.
It follows directly that
\begin{equation}\label{eq:gauge}
 B=s^{-1}PQ,\quad M_j=sPA_jP^{-1},\quad
 r=s^{-2}Pr_0,\quad C=sPSP^{-1}.
\end{equation}
Invertibility of $C$ implies invertibility of $S$. Substituting
\eqref{eq:gauge} in \eqref{eq:inverse} cancels the gauge matrices and
gives the scalar factor
$s^{-1-a+N-b-2}=1$, since $N=a+b+3$. The remaining expression is
$\alpha A_{w_1}\cdots A_{w_N}\beta$. Therefore
$F_{N,a}(\mu_a\psi)=\psi$ for every such parametrized tensor.

\emph{Closure and nonemptiness.}
Each $X_N$ is the closure of a polynomial image of an affine space,
so is irreducible. Equivalently, its vanishing ideal is the kernel of
the substitution homomorphism into the polynomial ring in the entries
of $A_0,A_1,\alpha,\beta$. This ring is an integral domain, and a
polynomial vanishing at all complex parameter values is zero. The
kernel, and hence the vanishing ideal, is prime.

The three-site tensor furnished by
\begin{equation}\label{eq:witness}
 A_0=\begin{pmatrix}1&2\\3&4\end{pmatrix},\quad
 A_1=\begin{pmatrix}5&6\\7&8\end{pmatrix},\quad
 l=(1,0),\quad r_0=(1,2)^T
\end{equation}
has lexicographically ordered coordinates
$(145,349,337,829,489,1173,1193,2933)$. Its five factors in
\eqref{eq:g} are respectively
\[
 387072,\quad172598092627968,\quad-397666005414838272,
 \quad-576,\quad2592.
\]
Thus $U$ is nonempty. Its inverse image under $\mu_a$ is nonempty by
the right-inverse construction. The parametrized tensors are dense in
$X_N$, and their intersection with this principal open subset is dense
in that subset. Both $F_{N,a}\mu_a$ and the identity are regular there,
so their coordinatewise equality extends from these parametrized tensors
to the whole subset. The two displayed maps are therefore inverse
isomorphisms of nonempty open subsets of the irreducible varieties.
\end{proof}

For explicit polynomial numerators and denominator, put $v=\det W$,
$c=\det J$. Formula~\eqref{eq:inverse} becomes
\begin{equation}\label{eq:polynomial}
 F_{N,a}(t)_w=
 \frac{\widetilde\ell(\adj J)^aN_{w_1}\cdots N_{w_N}
 (\adj J)^b\widetilde r}{v^3c^{N-3}}.
\end{equation}
Indeed $\ell=d^{-1}\widetilde\ell$,
$r=d^{-2}\widetilde r$, $M_j=(d/v)N_j$, and
$C^{-1}=(v/(dc))\adj J$; the powers of $d$ cancel because $a+b=N-3$.
The denominator is nonzero throughout $U$. This also makes the
regularity used in the proof explicit.

\section{Consequences}
\begin{corollary}\label{cor:rational}
For every $N\geq3$, $X_N$ is a rational variety of dimension seven.
\end{corollary}
\begin{proof}
In the normalized three-site representation~\eqref{eq:normalized}, write
$\ell=(x,1)$. The two identities $\ell M_i=e_i^T$ give
\[
 M_0=\begin{pmatrix}p&q\\1-xp&-xq\end{pmatrix},\qquad
 M_1=\begin{pmatrix}u&v\\-xu&1-xv\end{pmatrix}.
\]
The seven free parameters $x,p,q,u,v,r_0,r_1$ define a polynomial
map to $X_3$. Its image contains $U$ by the construction above.
Conversely, if a tensor from these seven parameters lies in $U$,
its matrix $P$ in the uniqueness argument is $I$ and $s=\ell_1=1$;
equation~\eqref{eq:gauge} returns the identical parameters. The
parameter preimage of $U$ is a nonempty open subset of $\C^7$,
and these maps identify it with $U$. Thus $X_3$ is birational to
$\C^7$; Theorem~\ref{thm:main} gives the assertion for $N\geq4$.
\end{proof}

\begin{corollary}\label{cor:weighted}
Let $3\leq m\leq N$, and retain any consecutive block of $m$ sites.
At each omitted site choose a fixed nonzero complex covector, allowing a
different choice at every site. Contracting the omitted sites with these
covectors defines a birational map $X_N\to X_m$.
\end{corollary}
\begin{proof}
First let $m=3$, with $a$ omitted sites on the left and $b$ on the right.
Write their ordered covectors as $\lambda^{(1)},\ldots,\lambda^{(a)}$
and $\eta^{(1)},\ldots,\eta^{(b)}$, respectively, and put
\[
 L(A)=\prod_{h=1}^a(\lambda^{(h)}_0A_0+\lambda^{(h)}_1A_1),\qquad
 R(A)=\prod_{h=1}^b(\eta^{(h)}_0A_0+\eta^{(h)}_1A_1),
\]
with products in site order and empty product $I$. Put $l=\alpha L(A)$ and
$r_0=R(A)\beta$. The contracted tensor has coordinates $t_{ijk}=lA_iA_jA_kr_0$
and belongs to $X_3$.
On the nonempty open subset of $U$ where $L(M)$ and $R(M)$ are invertible,
the inverse is
\begin{equation}\label{eq:weighted-inverse}
 t\longmapsto\bigl(\ell L(M)^{-1}M_{w_1}\cdots M_{w_N}R(M)^{-1}r\bigr)_w.
\end{equation}
The right-inverse identity follows by canceling $L(M),R(M)$.
For uniqueness apply the gauge argument~\eqref{eq:gauge} to these boundaries.
It gives
$L(M)=s^aPL(A)P^{-1}$ and $R(M)=s^bPR(A)P^{-1}$;
the same exponent $-1-a+N-b-2=0$ cancels all scalar factors.
Equality extends to the closures exactly as in Theorem~\ref{thm:main}.

To justify the additional open subset, each nonzero covector $\xi$ gives a
nonzero polynomial $\det(\xi_0A_0+\xi_1A_1)$ in the matrix parameters:
specialize $A_i=c_iI$ with $\xi_0c_0+\xi_1c_1\ne0$.
Their finite product and the nonzero chart polynomial have simultaneous
nonvanishing on the irreducible three-site parameter space. On that
intersection~\eqref{eq:gauge} transfers invertibility to $M$.
This proves nonemptiness for every fixed list of nonzero covectors.

For $m>3$, follow the $m$-site contraction $\nu$ by the ordinary contraction
$\sigma:X_m\to X_3$ retaining the first three sites of its block.
The composite $\sigma\nu$ is a three-site contraction of the type just proved,
with covector $(1,1)$ at the additional $m-3$ sites. Both $\sigma\nu$ and
$\sigma$ are birational; hence $\nu=\sigma^{-1}(\sigma\nu)$ is birational.
The case $m=N=3$ is the identity.
\end{proof}

\begin{corollary}\label{cor:minimal}
For $N\geq3$, no coherent contraction retaining at most two binary sites
is generically finite on $X_N$. Thus three is the least possible retained
block size for generic amplitude recovery of these models.
\end{corollary}
\begin{proof}
By Corollary~\ref{cor:rational}, $\dim X_N=7$; a tensor with at most two
binary indices lies in an affine space of dimension at most four.
The generic-fiber dimension theorem gives fiber dimension at least three.
This assertion concerns amplitudes and does not compare reduced density operators.
\end{proof}

\begin{corollary}\label{cor:equations}
For $N\geq4$ the $N-2$ quartics $\det K(\mu_a\psi)$ vanish on $X_N$,
but their common zero set is strictly larger than $X_N$.
Fix $a$, set $t=\mu_a\psi$, let $q=(\det W)^3(\det J)^{N-3}$,
and let $p_w$ be the numerator in~\eqref{eq:polynomial}.
On the ambient open subset $g(t)\ne0$, membership has the exact criterion
\begin{equation}\label{eq:membership}
 \psi\in X_N\quad\Longleftrightarrow\quad
 \det K(t)=0\ \text{ and }\ q(t)\psi_w=p_w(t)\quad\text{for all }w.
\end{equation}
Every member of $X_N$ on this chart is an actual tensor of the form
\eqref{eq:state}, not just a limit of such tensors.
\end{corollary}
\begin{proof}
Equation~\eqref{eq:outside} and the quartic relation prove the vanishing.
An affine zero locus defined by $N-2$ polynomials in $2^N$ variables has
dimension at least $2^N-(N-2)$ by the principal ideal theorem. This exceeds
seven for every $N\geq4$, so those quartics cannot define $X_N$ alone.
For~\eqref{eq:membership}, the forward implication is the inverse identity.
Conversely, Theorem 7 of~\cite{CM} identifies $X_3$ with $\det K=0$;
thus $t\in U$ and $q(t)\ne0$. The remaining equations give
$\psi=F_{N,a}(t)\in X_N$. Formula~\eqref{eq:inverse} supplies actual
matrix and boundary parameters for every such point.
\end{proof}

\section{Source coverage and verification}
The original posing source is the journal version of~\cite{CM}; its earlier
preprint is arXiv:1210.2812v1 (October 10, 2012), Section II, p.\ 4,
Conjecture 8. The same 22-term quartic already appears
in~\cite[Section 2, p.\ 4 of the arXiv version]{PS}.
The hidden-Markov birational result~\cite{HMM}
precedes that source. For the distinct periodic Conjecture 12, the proof in
\cite[Theorem 3.30]{CMS} was posted on April 16, 2019, and subsequently
published in 2023. The coverage below distinguishes these prior results from
the consequences proved here.

\begin{longtable}{@{}p{0.29\linewidth}p{0.65\linewidth}@{}}
\toprule
Source item in~\cite{CM} & Coverage\\
\midrule
Conjecture 8, Section 2, p.\ 6 & Proved: Theorem~\ref{thm:main}, every length and every consecutive
triple, on nonempty open subsets.\\
Theorem 7, p.\ 6 & Prior proof cited: the full quartic hypersurface description.
Proved here: $\det K$ vanishes and $X_3$ is rational of dimension seven
(Corollary~\ref{cor:rational}).\\
Paragraph before Conjecture 8, p.\ 6 & Proved: the $N-2$ quartic conditions;
they are not sufficient globally. Corollary~\ref{cor:equations} supplies exact
equations on the reconstruction chart.\\
Postselection comment, p.\ 6 & Proved: arbitrary nonzero site-dependent
coherent covectors and every retained block of length $m\geq3$
(Corollary~\ref{cor:weighted}); three amplitude indices are minimal
(Corollary~\ref{cor:minimal}).\\
Conjecture 9, pp.\ 6--7 & Not covered: open-boundary reduced-density
identifiability, including its minimality assertion.\\
Conjecture 10, p.\ 7 & Not covered: periodic reduced-density identifiability,
including its minimality assertion.\\
Conjecture 12, p.\ 9 & Prior proof cited:~\cite[Theorem 3.30, arXiv version]{CMS};
periodic trace-parameter identifiability, not Theorem~\ref{thm:main}.\\
Questions 1 and 2, p.\ 2 & The general periodic representability and defining-ideal
questions are not covered by the open-boundary reconstruction.\\
Theorems 3, 4; Propositions 5, 6, 11; Sections 1 and 3 & Prior proofs cited
in~\cite{CM}: periodic equations, homogeneity, symmetry and parametrization;
not new consequences of the open-boundary argument.\\
Section 3 HMM comparison; Section 4 extensions & Prior proof cited:
\cite[Theorems 4.1, 5.11]{HMM} for the normalized HMM models.
Larger bond/physical dimensions and other network types are not covered.\\
\bottomrule
\end{longtable}

\paragraph{Lean and finite checks.}
In the
\href{https://github.com/the-omega-institute/trureturing/blob/b95cdbb8cccca425964251ba38cc0b0e517e478a/D5/S3/Quantum/TensorNetworks/OpenBoundaryMPSMarginalBirationality.lean}{pinned Lean file},
the declaration \texttt{result} in the namespace
\texttt{D5.S3.Quantum.TensorNetworks.}\allowbreak
\texttt{OpenBoundaryMPSMarginalBirationality}
has type \texttt{claim}. Its coverage is the two prime vanishing ideals and
\texttt{IsBirationalVia}: polynomial $g,q,p_w$, a nonempty chart, a nonzero
denominator and both inverse identities for every $N\geq4$, $a+3\leq N$.
Equations~\eqref{eq:g} and~\eqref{eq:polynomial} match that chart and those
rational coordinates. Corollaries~\ref{cor:rational}--\ref{cor:equations}
are written consequences, not additional declarations of that file.
The accompanying \texttt{verify.py} compares all coefficients of
$\det K$ with the 22-term quartic in~\cite[Theorem 7]{CM}, checks the witness,
and evaluates both inverse identities exactly for six rational parameter
choices, all lengths $4$ through $8$, and every consecutive triple.
Its 120 source cases, 120 target cases and 47,040 coordinate equalities
are finite checks; the proof above supplies the all-length conclusion.
The separate \texttt{verify\_extensions.py} checks ordered site-dependent
contractions and longer-block reconstruction with exact rational arithmetic;
its sample ranges and counts appear in \texttt{extension-checks.json}.

\begin{thebibliography}{9}
\bibitem{CM} A. Critch and J. Morton,
\emph{Algebraic geometry of matrix product states},
SIGMA \textbf{10} (2014), 095, 10 pp.
\href{https://doi.org/10.3842/SIGMA.2014.095}{doi:10.3842/SIGMA.2014.095}.
\bibitem{HMM} A. Critch,
\emph{Binary hidden Markov models and varieties},
J. Algebraic Statistics \textbf{4} (2013), 1--30.
\href{https://arxiv.org/abs/1206.0500v3}{arXiv:1206.0500v3}.
\bibitem{CMS} A. Czapli\'nski, M. Micha\l{}ek and T. Seynnaeve,
\emph{Uniform matrix product states from an algebraic geometer's point of view},
Adv. Appl. Math. \textbf{142} (2023), article 102417.
\href{https://doi.org/10.1016/j.aam.2022.102417}{doi:10.1016/j.aam.2022.102417}.
Theorem numbering here follows \href{https://arxiv.org/abs/1904.07563v1}{arXiv:1904.07563v1}.
\bibitem{FLAC} M. Florido-Llin\`as, \'A. M. Alhambra,
D. P\'erez-Garc\'ia and J. I. Cirac,
\emph{Uniform matrix product states with a boundary},
\href{https://arxiv.org/abs/2512.11968v1}{arXiv:2512.11968v1} (2025).
\bibitem{LS} C. De Lazzari, H. J. Motwani and T. Seynnaeve,
\emph{The linear span of uniform matrix product states},
SIGMA \textbf{18} (2022), 099, 18 pp.
\href{https://doi.org/10.3842/SIGMA.2022.099}{doi:10.3842/SIGMA.2022.099}.
\bibitem{PS} L. Pachter and B. Sturmfels,
\emph{Tropical geometry of statistical models},
Proc. Natl. Acad. Sci. USA \textbf{101} (2004), 16132--16137.
\href{https://doi.org/10.1073/pnas.0406010101}{doi:10.1073/pnas.0406010101}.
The quartic locator follows \href{https://arxiv.org/abs/q-bio/0311009v2}{arXiv:q-bio/0311009v2}.
\end{thebibliography}
\medskip
\noindent\small\textcopyright\ 2026 Haobo Ma and Wenlin Zhang.
Licensed under \href{https://creativecommons.org/licenses/by/4.0/}{Creative Commons Attribution 4.0 International (CC BY 4.0)}.
\end{document}
